\documentclass[11pt]{article}
\usepackage[margin=1in]{geometry}
\usepackage{booktabs}
\usepackage{tabularx}
\usepackage{array}
\usepackage{enumitem}
\usepackage{hyperref}
\usepackage[T1]{fontenc}
\usepackage{lmodern}

\newcolumntype{Y}{>{\raggedright\arraybackslash}X}
\setlength{\parindent}{0pt}
\setlength{\parskip}{6pt}

\title{DS-GA 1011 Assignment 4}
\author{Boyu Chen}
\date{}

\begin{document}
\maketitle

\textbf{Code (GitHub):} \url{https://github.com/boyuchencoding/dsga1011-a4}\\
\textbf{Q7 model checkpoint (Google Drive):} \url{https://drive.google.com/drive/folders/13cO9wzcfDiukEQS7DSnA-_ljzPwiTYPX?usp=sharing}

All experiments were run locally on an NVIDIA RTX 4070 Laptop GPU (8\,GB).

\section{Part 1: Fine-tuning BERT for Sentiment Classification}

\subsection{Q1. Fine-tuning BERT (\texttt{bert\_finetuning})}

The training loop is in \texttt{part1/main.py} (\texttt{do\_train}). For every epoch and every batch, the
batch is moved to the GPU, a forward pass with the labels returns the cross-entropy loss, and then
\texttt{loss.backward()}, \texttt{optimizer.step()}, \texttt{lr\_scheduler.step()} and
\texttt{optimizer.zero\_grad()} are called. The default hyperparameters were kept
(\texttt{bert-base-cased}, AdamW, learning rate $5\times10^{-5}$, linear decay, batch size 8, 3 epochs).

\begin{table}[h]
\centering
\begin{tabular}{lr}
\toprule
Run & Accuracy \\
\midrule
Debug (4k train / 1k test subset) & 89.70 \\
Full training, original test set (25k) & \textbf{93.16} \\
\bottomrule
\end{tabular}
\caption{BERT fine-tuning results. Output file: \texttt{out\_original.txt}.}
\end{table}

\subsection{Q2a. Data transformation (\texttt{data\_transformation\_design})}

\textbf{Transformation: keyboard typos.} Each review is split on whitespace and on HTML tags
(e.g.\ \texttt{<br />}); separators are kept so the text can be rebuilt exactly. Every token that is purely
alphabetic ASCII and has at least 4 letters is selected independently with probability $p=0.5$. In a selected
word, one letter position is chosen uniformly at random and that letter is replaced by a uniformly random
neighbouring key on a QWERTY keyboard (e.g.\ \texttt{a}$\to$\{\texttt{q,s,z}\}, \texttt{o}$\to$\{\texttt{i,p,l}\},
\texttt{m}$\to$\{\texttt{n,j,k}\}); the case of the letter is preserved. Short words ($<4$ letters), numbers,
punctuation and HTML tags are never changed. The random generator is seeded (\texttt{random.seed(0)}), so the
transformed test set is reproducible.

\textit{Example:} ``When I unsuspectedly rented A Thousand Acres, \ldots\ of course Michelle Pfeiffer was in it''
$\to$ ``When I unsuspectedly rented A Thousan\textbf{c} Acres, \ldots\ of c\textbf{p}urse M\textbf{k}chelle
Pfeiffer was in it''.

\textbf{Why it is reasonable.} Hitting an adjacent key is the most common kind of typo when people type
quickly, especially on phones, and online reviews are full of such mistakes. Every changed word differs by a
single character, so a human can still read the review and would give it the same sentiment label. At the
same time the transformation is not trivial for BERT: a cased WordPiece tokenizer splits a misspelled word
(``wonderfuoly'') into rare sub-word pieces that the model rarely saw in training, so the sentiment cue carried
by that word can be lost.

\subsection{Q2b. Implementation (\texttt{data\_transformation\_impl})}

Implemented in \texttt{custom\_transform} in \texttt{part1/utils.py} (constants \texttt{TYPO\_PROB=0.5} and
\texttt{QWERTY\_NEIGHBORS}). I first tried $p=0.3$, which only dropped accuracy by 2.8 points, and increased it to
$p=0.5$.

\begin{table}[h]
\centering
\begin{tabular}{lrr}
\toprule
Test set & Accuracy & Change \\
\midrule
Original & 93.16 & -- \\
Transformed, $p=0.3$ (not used) & 90.34 & $-2.82$ \\
Transformed, $p=0.5$ (final) & \textbf{87.95} & $\mathbf{-5.21}$ \\
\bottomrule
\end{tabular}
\caption{Model from Q1 on the original and transformed test sets. Output file: \texttt{out\_transformed.txt}.}
\end{table}

\subsection{Q3. Data augmentation (\texttt{data\_augmentation})}

\texttt{create\_augmented\_dataloader} (\texttt{part1/main.py}) samples 5{,}000 random training examples
(\texttt{shuffle(seed=42).select(range(5000))}), applies \texttt{custom\_transform} to them, concatenates them
with the full 25k original training set (30k examples in total, 15{,}006 positive / 14{,}994 negative),
tokenizes them exactly like the original data and returns a shuffled \texttt{DataLoader}.

\textbf{Note on the learning rate.} With the default learning rate $5\times10^{-5}$ the augmented run diverged:
the model predicted the same class for every input (50.0\% on both test sets). I checked that the augmented
dataset itself was correct (sizes, labels and transformed text), so this was a training instability, and I
retrained with learning rate $2\times10^{-5}$ (all other settings unchanged).

\begin{table}[h]
\centering
\begin{tabular}{lrr}
\toprule
Model & Original test & Transformed test \\
\midrule
Q1 model (original data, lr $5\times10^{-5}$) & 93.16 & 87.95 \\
Augmented, lr $5\times10^{-5}$ (diverged) & 50.00 & 50.00 \\
Augmented, lr $2\times10^{-5}$ (final) & \textbf{93.52} & \textbf{90.60} \\
\bottomrule
\end{tabular}
\caption{Effect of data augmentation. Output files: \texttt{out\_augmented\_original.txt},
\texttt{out\_augmented\_transformed.txt}.}
\end{table}

\textbf{(1) Transformed test set.} Accuracy improved from 87.95 to 90.60 ($+2.65$), so the gap to the original
test set shrank from 5.2 to 2.9 points. Augmentation recovers about half of the drop, but not all of it.

\textbf{(2) Original test set.} Accuracy did not get worse; it went from 93.16 to 93.52 ($+0.36$), which is
within run-to-run noise. Augmentation did not hurt clean performance.

\textbf{Intuitive explanation.} The Q1 model never saw misspelled words, so the unusual sub-word pieces that typos
create carry no learned sentiment signal. In the augmented data, 5k reviews contain these pieces with the same
labels as before, so the model learns that ``wonderfuoly'' behaves like ``wonderfully'' and relies less on the
exact spelling of a few key words. Because the augmented examples keep their labels and the original 25k
examples are all still there, the clean distribution is not distorted and accuracy on the original test set
stays the same. One caveat: the final augmented model used a smaller learning rate than the Q1 model, so a
small part of the difference may come from the learning rate rather than from augmentation.

\textbf{Limitation.} The augmentation only covers the exact shift we designed (QWERTY single-letter typos with a
fixed rate). Real OOD data can shift in many other ways (synonyms, slang, abbreviations, other domains, other
typo patterns), and the model is not made robust to those. To handle each new shift we would have to know it in
advance, write a new transformation and retrain, which does not scale and cannot cover unknown shifts. Even for
the known shift the gain is partial (2.9 points remain), since 5k transformed examples are only 17\% of the
training data.

\section{Part 2: Fine-tuning T5 for Text-to-SQL}

%%PART2%%

\end{document}
